%%
%% This is file `tikzposter-example.tex',
%% generated with the docstrip utility.
%%
%% The original source files were:
%%
%% tikzposter.dtx  (with options: `tikzposter-example.tex')
%% 
%% This is a generated file.
%% 
%% Copyright (C) 2014 by Pascal Richter, Elena Botoeva, Richard Barnard, and Dirk Surmann
%% 
%% This file may be distributed and/or modified under the
%% conditions of the LaTeX Project Public License, either
%% version 2.0 of this license or (at your option) any later
%% version. The latest version of this license is in:
%% 
%% http://www.latex-project.org/lppl.txt
%% 
%% and version 2.0 or later is part of all distributions of
%% LaTeX version 2013/12/01 or later.
%% 








 \documentclass[17pt, a1paper, portrait, margin=0mm, innermargin=10mm,
     blockverticalspace=15mm, colspace=10mm, subcolspace=8mm]{tikzposter} %Default values for poster format options.

 %\tikzposterlatexaffectionproofon %shows small comment on how the poster was made at bottom of poster
\usepackage{exscale}
\usepackage{listings}
\usepackage{hyperref}
\usepackage{wrapfig}

 % Commands
 \newcommand{\bs}{\textbackslash}   % backslash
 \newcommand{\cmd}[1]{{\bf \color{red}#1}}   % highlights command

 % Title, Author, Institute
\title{Can Twitter User's Moods Predict the Stock Market?}
\author{Aaron Gonzales and Adam Delora}
\institute{Department of Computer Science, University of New Mexico} 
%\titlegraphic{blah}

 % -- PREDEFINED THEMES ---------------------- %
 % Choose LAYOUT:  Default, Basic, Rays, Simple, Envelope, Wave, Board, Autumn, Desert,
 \usetheme{Rays}
%\usecolorstyle[colorPalette=BrownBlueOrange]{Germany}


 \begin{document}

 \maketitle[width=55cm, titletotopverticalspace=0.4cm]

     \begin{columns}%blocks will be placed into columns
         \column{.50}
%%%%%%%%%%%%%%%%new block%%%%%%%%%%%%%%%%%%%%%%%%%%%%%%%%%%%%
      \block{Background}{
          \begin{itemize}
          \item Twitter is a microblogging service with 284 million monthly
            active users who post 500 million updates (``tweets'') per day.
          \item Latent Semantic Indexing is a technique used to summarize
            words (documents) into representative ideas similar to primary
            component analysis.
          \item The AFINN semantic indexing database assigns coded valence
            values to common words to allow quantification of a set of word's
            ``mood''.
        \end{itemize}
          }
%%%%%%%%%%%%%%%%new block%%%%%%%%%%%%%%%%%%%%%%%%%%%%%%%%%%%%
       \block{Methods}{
          We collected tweets using the public Twitter Streaming RESTful API
          over 2014-10-15 \textemdash 2014-11-09 by tracking words that related
          to various tech stocks indexed by NASDAQ.           
          \begin{tikzfigure}[Several examples of tweets] \label{fig:tweets}
            \begin{itemize}
              \item \small{Maine nurse defies Ebola quarantine order by taking bike ride.
                http://t.co/eERINkm3AQ via @indystarJ}
              \item \small{Tim Cook: Apple CEO Says 'Being Gay Is One Of God's
                Greatest Gifts' Earlier today, the chief executive of
                Apple,\ldots \url{http://t.co/zIXb2HbDmd}}
              \item \small{ Meet the Swedish twin sisters who want to be
                'identical artificial dolls' \url{http://t.co/77eA3esaQa}}
              \item \small{Last Christmas I got a black iPhone it was the
                worst Christmas present ever, I wanted a white one. I mean I'm
                not\ldots \url{http://t.co/jhP9ODs3QK}}
            \end{itemize}
            \tiny{\textcolor{white}{blah}}
          \end{tikzfigure}
          %}
          %\begin{wrapfigure}[6]{l}{0.33\linewidth}
          \begin{tikzfigure}[Kewords tracked] \label{fig:keywords}
            \begin{tikzpicture}
              \node[anchor=north west] (note2) at (0,0)
              { \small{ebola, aapl, apple, mac, tim ook, goog,
                  google, gmail, youtube, microsoft, msft,
                  nadella, twrt, amazon, amzn, prime, aws, fb,
                  facebook}};
            \end{tikzpicture}
          \end{tikzfigure}
        %\end{wrapfigure}
          NASDAQ market data was also collected over the same period. Tweets
          were preprocessed to remove common stopwords and latent semantic
          indexing was performed on one-hour bag-of-words bins of tweets,
          giving a total of xxx hours included in analysis. Each hour bin's LSI
          topics were scored using the AFINN database, resulting in a single
          number indicating semantic valence for each hour. The LSI score was
          smoothed using rolling means (Figure \ref{fig:rollingmeans} and
          assessed for periodicity.  semantic data was combined with the stock
          data for visualization and analysis. A vector autoregressive (VAR)
          model was performed to assess the predictive power of the semantic
          data against the stock data.

       Using LSI would provide something like the following topic:
       \begin{tikzfigure}[Two example LSI topics] \label{fig:lsitopic}
         \small{
           ``0.748*+ 0.199*``video'' + 0.190*``my'' + -0.183*``i''
           + -0.172*``yo'' + 0.153*``me'' + 0.147*``music'' +
           -0.140*``apple'' + 0.114*``give'' + 0.113*``be'''

           '0.416*``scientists'' + 0.219*``are'' + 0.217*``about'' +
           0.216*``do'' + 0.214*``don\'t'' + 0.211*``over'' + 0.209*``say'' +
           0.209*``climate'' + 0.209*``it!'' + 0.208*``panic'''

           '0.586*``on'' + -0.315*``is'' + -0.217*``get'' + 0.165*``from''
           + -0.140*``ebola'' + -0.129*``i'' + -0.126*``google'' +
         -0.122*``it'' + -0.117*``really'' + 0.114*``liked'''}
            %\caption{Two exaple LSI topics}
       \end{tikzfigure}
        } 
%%%%%%%%%%%%%%%%new block%%%%%%%%%%%%%%%%%%%%%%%%%%%%%%%%%%%%
     \block{Summary Statistics}{
       \begin{wrapfigure}[10]{r}{0.35\linewidth}
         \begin{tikzfigure}[Rolling sample.] \label{fig:summary}
             %\includegraphics[width=4in]{figures/lsi_rolling_means.png}
           \includegraphics[width=9cm]{figures/lsi_rolling_means.png}
         \end{tikzfigure}
       \end{wrapfigure}
       Something results are plotted in Figure \ref{fig:summary}
       \begin{itemize}
         \item Total tweets collected: 86 million
         \item Total number of days: 25
         \item Total hour bins: 360
       \end{itemize}

       \begin{tikzfigure}[ test ]
       \begin{tikzpicture}
         \node[anchor=south west,inner sep=0] (image) at (0,0)
         {\includegraphics[width=0.9\linewidth]{figures/placeholder}};
         \begin{scope}[x={(image.south east)},y={(image.north west)}]
           \draw[red,ultra thick,rounded corners] (0.62,0.65)
           rectangle (0.78,0.75);
         \end{scope}
       \end{tikzpicture}
     \end{tikzfigure}
     }

%%%%%%%%%%%%%%%%%%%%%%%%%%%%%%%%%%%%%%%%%%%%%%%%%%%%%%%%%%
%%%%%%%%%       new col         %%%%%%%%%%%%%%%%%%%%%%%%%%
%%%%%%%%%%%%%%%%%%%%%%%%%%%%%%%%%%%%%%%%%%%%%%%%%%%%%%%%%%

     \column{.5}
%%%%%%%%%%%%%%%%new block%%%%%%%%%%%%%%%%%%%%%%%%%%%%%%%%%%%%
     \block{LSI}{
       Using LSI would provide something like the following topic:
       \begin{tikzfigure}[Rolling means of LSI score per hour over
         sample.] \label{fig:means}
         \includegraphics[width=7in]{figures/lsi_rolling_means.png}
       \end{tikzfigure}
     }
%%%%%%%%%%%%%%%%new block%%%%%%%%%%%%%%%%%%%%%%%%%%%%%%%%%%%%
     \block{Columns}{
       By default, blocks are arranged in a single column. If you want
       multiple columns for your poster, you may use the
       \texttt{columns} environment. For example, 
       \begin{quote}
         \texttt{\noindent \bs begin\{columns\}\\
         \bs column\{.6\}\\
         \bs block\{\dots\}\{\dots\}\\
         \bs column\{.4\}\\
         \bs block\{\dots\}\{\dots\}\\
         \bs block\{\dots\}\{\dots\}\\
         \bs end\{columns\}
       }
     \end{quote}
     will create two columns of 60\% and 40\% the available width;
     spacing between successive columns is handled automatically.  The
     block command(s) following \texttt{\bs column} are the blocks to
     go in that column.  The number of columns is free to be chosen, but
     the relative widths must all be chosen.  If the widths sum to less
     than 1, empty space will be seen on the right.  If they sum to more than
   1, the latter columns will be cut off.  }

         \begin{subcolumns}
             \subcolumn{.45}
             \block{Subcolumns}{If you want to have an additional subdivision
               of columns inside a column, you may use the\\ \texttt{\bs
               subcolumns} environment inside of a column environment.  The
             functionality is similar to that of columns, but now the widths
           are relative to the width of the current column.}

             \subcolumn{.5}
             \block{}{An example use of subcolumns is.
                 \begin{quote}
                     \texttt{\bs begin\{subcolumns\}\\
                     \bs subcolumn\{.6\}\\
                     \bs block\{\dots\}\{\dots\}\\
                     \bs subcolumn\{.4\}\\
                     \bs block\{\dots\}\{\dots\}\\
                     \bs block\{\dots\}\{\dots\}\\
                     \bs end\{subcolumns\}
                     }
             \end{quote}
         }
         \end{subcolumns}

         \block[titlewidthscale=.8,bodywidthscale=.9,titleoffsety=9.5mm,bodyoffsety=9mm]{Changing the Poster's Appearance}{
             If the default appearance of the title, background, blocks, and
             notes is not desired, you may change the colors by calling the
             color style along with a general layout theme with the commands
             \begin{quote}
       \texttt{\bs usecolorpalette}\{{\em color palette}\}\\
                 \texttt{\bs usecolorstyle\{{\em color style}\}}
             \end{quote}
             and
             \begin{quote}
                 \texttt{\bs usetheme\{{\em layout style}\}}
             \end{quote}
             where the color palette and style and layout style are either the name of a custom made or one of the offered predefined choices listed in the manual or the comments of this poster's source.  Individual changes can be made to the style of the  background, title matter, blocks, inner blocks, and notes by using one of the following (along with either a custom-designed style or a predefined style listed in the manual or the comments of this poster's source).  These changes are made with the commands
             \begin{quote}
                 \texttt{\bs usebackgroundstyle[]\{\}, \bs usetitlestyle[]\{\},\\ \bs useblockstyle[]\{\},\bs innerblockstyle[]\{\}, \bs usenotestyle[]\{\}}
             \end{quote}
             Custom styles for these can be made; this is detailed in the manual.
          }

     \end{columns}

     \block[titleoffsety=-1cm,bodyoffsety=-1cm]{Sample document}{\vspace{2em}
         This poster was created by the following commands (omitting the contents of the blocks and notes) to give a sense of how different objects are created and options used.
         \begin{quote}
             \texttt{\bs documentclass[25pt, a0paper, portrait, margin=0mm, innermargin=15mm,
         blockverticalspace=15mm, colspace=15mm, subcolspace=8mm]\{tikzposter\}\\
             \bs title\{Using tikzposter\} \bs author\{Pascal Richter, Elena Botoeva, Richard Barnard, \& Dirk Surmann\} \bs institute\{\}\\
              \bs usetheme\{Autumn\}\bs usecolorstyle[colorPalette=BrownBlueOrange]\{Germany\}\\
             \bs begin\{document\}\bs maketitle\\
             \bs begin\{columns\} \bs column\{0.55\}\\
             \bs block\{Creating the document\}\{The document\dots\} \bs note[targetoffsetx=-.05\bs textwidth,targetoffsety=9.5cm,innersep=.4cm,angle=-45,connection]\{\dots\}\\
             \bs block\{The title matter\}\{The title\dots\}\\
             \bs block\{Blocks\}\{Blocks are\dots\} \bs note[targetoffsetx=24cm, targetoffsety=-9cm,rotate=1,angle=270,radius=8cm,width=.75\bs textwidth,innersep=.4cm]\{You can\dots\}\\
             \bs column\{0.45\} \bs block\{Columns\}\{By default,\dots\}\\
             \bs begin\{subcolumns\} \bs subcolumn\{.45\}
             \bs block\{Subcolumns\}\{If you\dots\}
             \bs subcolumn\{.5\} \bs block\{\}\{An example\dots\}
             \bs end\{subcolumns\}\\
             \bs block[titlewidthscale=.8,bodywidthscale=.9,titleoffsety=9.5mm,bodyoffsety=9mm]\{Changing the Poster's Appearance\}\{If the default\dots\}
             \bs end\{columns\}\\
             \bs block[titleoffsety=-1cm,bodyoffsety=-1cm]\{Sample document\}\{This poster\dots\}\\
             \bs end\{document\}
             }
         \end{quote}
     }

 \end{document}




\endinput
%%
%% End of file `tikzposter-example.tex'.
